% ~~~~~~~~~~~~~~~~~~~~~~~~~~~~~~~~~~~~~~~~~~~~~~~~~~~~~~~~~~~~~~~~~~~~~~ %
% ~~~~~~~~~~~~~~~~~~~~~ Projektspezifische Makros ~~~~~~~~~~~~~~~~~~~~~~ %
% ~~~~~~~~~~~~~~~~~~~~~~~~~~~~~~~~~~~~~~~~~~~~~~~~~~~~~~~~~~~~~~~~~~~~~~ %

% Shortcuts für häufig wiederkehrende Begriffe
\newcommand{\asw}{ASW}

% Wenn nach den 3 Ebenen nach \subsubsection eine weitere Gliederungsebene benötigt wird
\newcommand{\paragraphheader}[1]{\paragraph{#1}\mbox{}\\}

% Vordefinierte Arten der Fußnoten. 
% Die Struktur kann an dieser Stelle geändert werden und betrifft jede Verwendung im gesamten Dokument,
% was praktisch sein kann, falls der Betreuer oder Gutachter zb. die ausgeschriebene Variante von "Vgl." bevorzugt.
\newcommand{\vgl}[2]{\footcite[Vgl.][#2]{#1}} % chktex 9 chktex 10
\newcommand{\footrefnote}[2]{\footnote{Für das Thema #1 siehe Kapitel~\ref{#2}}}
\newcommand{\wholesection}[2]{\footcite[Für den gesamten Abschnitt vgl.][#2]{#1}} % chktex 9 chktex 10

% ~~~~~~~~~~~~~~~~~~~~~~~~~~~~~~~~~~~~~~~~~~~~~~~~~~~~~~~~~~~~~~~~~~~~~~ %
% ~~~~~~~~~~~~~~~~~~~~~~~~~~~ Silbentrennung ~~~~~~~~~~~~~~~~~~~~~~~~~~~ %
% ~~~~~~~~~~~~~~~~~~~~~~~~~~~~~~~~~~~~~~~~~~~~~~~~~~~~~~~~~~~~~~~~~~~~~~ %

% Falls LaTeX Wörter nicht richtig trennt können hier eigene Angaben in folgendem Stil gemacht werden
\hyphenation{Java-Script} % chktex 17
